% ============================================================================
%  Debt-Financed Fiscal Stimulus, Heterogeneity, and Welfare
%  MODEL DESCRIPTION WITH ENDOGENOUS INFLATION  +  COUNTERFACTUAL RESULTS
%
%  Standalone document: compiles on its own, no .bib required.
%  To move into Overleaf, paste the body between \begin{document} and
%  \end{document} into the manuscript; biblatex there will take over \textcite
%  and the \defcite shim in this preamble becomes inert.
%
%  All numbers in Sections 5-6 are generated by
%  scripts/table6_consolidate_sw.jl and live in output/table6_sw.md/.txt.
% ============================================================================
\documentclass[11pt]{article}

\usepackage[margin=1.1in]{geometry}
\usepackage{amsmath,amssymb,amsthm}
\usepackage{dsfont}
\usepackage{booktabs}
\usepackage{framed}
\usepackage{xcolor}
\usepackage{graphicx}
\usepackage{caption}
\usepackage[hidelinks]{hyperref}

\definecolor{shadecolor}{gray}{0.94}
\newcommand{\Ind}{\mathds{1}}

% ---- citation shim: makes the file compile without a .bib -------------------
\newcommand{\citeLookup}[1]{%
  \ifcsname CTX#1\endcsname\csname CTX#1\endcsname\else\textsf{[#1]}\fi}
\providecommand{\textcite}[1]{\citeLookup{#1}}
\newcommand{\defcite}[2]{\expandafter\def\csname CTX#1\endcsname{#2}}
\defcite{Bewley-77}{Bewley (1977)}
\defcite{Imrohoroglu-89}{\.{I}mrohoro\u{g}lu (1989)}
\defcite{Aiyagari-94}{Aiyagari (1994)}
\defcite{Huggett-1996}{Huggett (1996)}
\defcite{Epstein-Zin-89}{Epstein and Zin (1989)}
\defcite{Greenwood_Hercowitz_Huffman_1988}{Greenwood, Hercowitz and Huffman (1988)}
\defcite{Heathcote-Storesletten-Violante-2017}{Heathcote, Storesletten and Violante (2017)}
\defcite{Hatchondo-Martinez-09}{Hatchondo and Mart\'{\i}nez (2009)}
\defcite{Fourakis-2024}{Fourakis (2024)}
\defcite{Rotemberg-1982}{Rotemberg (1982)}
\defcite{Erceg-Henderson-Levin-2000}{Erceg, Henderson and Levin (2000)}
\defcite{Auclert-Rognlie-Straub-2024}{Auclert, Rognlie and Straub (2024)}
\defcite{Taylor-1993}{Taylor (1993)}
\defcite{Schmitt-Grohe-Uribe-2005}{Schmitt-Groh\'e and Uribe (2005)}

\title{\textbf{Debt-Financed Fiscal Stimulus, Heterogeneity, and Welfare}\\[0.4em]
\large Model with Endogenous Inflation, and Counterfactual Results}
\author{Leonardo Barreto}
\date{\today}

\begin{document}
\maketitle

\begin{abstract}
\noindent This document restates the model of the paper with the exogenous
inflation assumption removed. Inflation is now determined inside the model by a
sticky-wage block: a continuum of labor unions set nominal wages subject to a
quadratic adjustment cost, households supply hours under proportional rationing,
and a monetary authority follows a Taylor rule. Sections
\ref{model}--\ref{calibration} are the model, equilibrium and calibration;
Sections \ref{results}--\ref{decomposition} report the counterfactual results
under the extension, together with a decomposition of the one channel that the
extension adds and the flexible-wage model does not have.
\end{abstract}

% ============================================================================
\section{Model} \label{model}
% ============================================================================

I build an open economy heterogeneous agent model in the spirit of
\textcite{Bewley-77}, \textcite{Imrohoroglu-89}, \textcite{Aiyagari-94}, and
\textcite{Huggett-1996}. Time is discrete and indexed by $t\in\{1,2,\dots\}$.
There is no aggregate uncertainty. \textbf{Inflation is determined
endogenously}: nominal wages are sticky, so the price level is pinned down
jointly by a wage Phillips curve, a monetary policy rule and the Fisher
equation. The economy is populated by domestic and foreign households, final
good firms, a continuum of labor unions, a monetary authority, and a government.
Domestic households consume final goods, supply labor at hours set by their
union, face time-varying idiosyncratic income risk, and self-insure by saving
and borrowing in capital and long-term nominal bonds. Foreign households also
demand domestic bonds. Firms produce the final good using a
constant-returns-to-scale technology with capital and labor as inputs. The
government pays out transfers to households and finances its expenditures
through labor and capital income taxes as well as bond issuance. I consider a
perfect foresight transition sequence between an initial stationary equilibrium
with a given stock of government debt and a new stationary equilibrium with a
higher debt-to-GDP ratio following a one-time zero-probability shock. In the
following subsections, I describe each block of the model.

\subsection{Domestic Households} \label{households_problem}

\paragraph{Preferences.} The economy is populated by a unit continuum of
ex-ante identical households indexed by $i\in[0,1]$. Household preferences are
given by the following \textcite{Epstein-Zin-89} utility function:
\begin{equation}\label{utility_function}
V_{it}= \left\{(1-\beta)x_{it}^{1-1/\xi} + \beta \left[\mathbb{E}\left( V_{it+1}^{1-\gamma}\right) \right]^{\frac{1-1/\xi}{1-\gamma}} \right\}^{\frac{1}{1-1/\xi}},
\end{equation}
where $\mathbb{E}$ denotes the expectation operator over idiosyncratic states,
$\beta\in(0,1)$ is the discount factor, $\xi$ is the elasticity of
intertemporal substitution, $\gamma$ is the risk aversion coefficient, and $x$
is the \textcite{Greenwood_Hercowitz_Huffman_1988} composite of consumption and
labor:
\begin{equation}\label{composite_c_n}
    x_{it} = c_{it} - \theta \frac{n_{it}^{1+1/\nu}}{1+1/\nu},
\end{equation}
where $c_{it}$ denotes consumption of the final good in period $t$, which is
sold at a nominal price $P_t$, and $n_{it}$ denotes the amount of labor supplied
in period $t$. The parameter $\nu$ controls the Frisch elasticity of labor
supply, and $\theta$ is a labor disutility parameter. The GHH composite plays a
larger role here than in the flexible-wage version of the model: because the
marginal disutility of work is separable from consumption, the steady-state
hours profile is independent of assets, which is what makes the proportional
rationing scheme of Section \ref{unions} internally consistent.

\paragraph{Idiosyncratic income risk.} Households face uninsurable risk to
their individual incomes. Let $z_{it}\geq 0$ be the idiosyncratic component of
household $i$'s income at date $t$. This idiosyncratic component follows a
Markov process with ergodic mean normalized to one, $\mathbb{E}(z_{it})=1$ for
all $t$, evolves according to a time-varying transition matrix $\Omega_t$, and
has an ``unemployment state'' in which $z_{it}=0$. The probability that an
unemployed worker stays unemployed, denoted as the U-U transition probability,
is $P^{uu}_t \equiv \text{Prob}(z_{it}=0\mid z_{it-1}=0)$, and the probability
that an employed worker remains employed, referred to as the E-E transition
probability, is $P^{ee}_t \equiv \text{Prob}(z_{it}>0\mid z_{it-1}>0)$. These
probabilities help capture the dynamics of the labor market since 2020.

\paragraph{Hours under sticky wages.} This is the first point at which the
model departs from the flexible-wage version. Nominal wages are set by unions,
not by households, and hours are rationed \emph{proportionally}: every employed
household works the same fraction of its own steady-state hours. Formally, let
$\tilde n(z)$ denote the hours a household with productivity $z$ works in the
initial stationary equilibrium, and let $s_t$ be a common, economy-wide
\emph{rationing index}. Then
\begin{equation}\label{rationing}
  n_{it} \;=\; s_t\,\tilde n(z_{it}),
\end{equation}
with $\tilde n(0)=0$, so unemployed households supply no labor. Households take
$n_{it}$ as given and choose only consumption and savings. Aggregate effective
labor is therefore
\begin{equation}\label{agg_labor_rationed}
  N_t \;=\; s_t \sum_{z} \Lambda_t(z)\, z\, \tilde n(z),
\end{equation}
where $\Lambda_t(z)$ is the date-$t$ marginal distribution of productivity.

Proportional rationing is not an innocuous normalization, and it is worth being
explicit about why it is the right one here. Under flexible wages and the GHH
composite \eqref{composite_c_n} with the tax scheme \eqref{tax_system}, the
household's intratemporal condition delivers hours that are proportional to
$\tilde n(z)$ with a factor that is \emph{common across $z$}: the productivity
term $z^{1-\tau}$ cancels exactly. That common factor is
\begin{equation}\label{s_flex}
  s^{\text{flex}}_t \;=\;
  \left[\frac{\tau_t^n}{\tau^n_{ss}}\left(\frac{w_t}{w_{ss}}\right)^{1-\tau}\right]^{\frac{\nu}{1+\tau\nu}} ,
\end{equation}
so the family of hours profiles spanned by \eqref{rationing} contains every
flexible-wage allocation. Setting $s_t = s^{\text{flex}}_t$ recovers the
flexible-wage model exactly, which nests the original paper as a limiting case
and provides a direct numerical check on the extension.

\paragraph{Budget and borrowing constraints.} When households provide
$n_{it}\geq0$ units of labor, they receive $W_tz_{it}n_{it}$ in labor earnings,
where $W_t$ is the nominal wage rate per efficiency unit of labor. Labor income
is taxed progressively according to a
\textcite{Heathcote-Storesletten-Violante-2017} scheme $\mathcal{T}$ that
depends on a rate $\tau^n_t$, so high-income households pay labor income taxes
while low-income households receive a net transfer. Households also receive
unemployment benefits $P_tT^u_t$ when unemployed, and lump-sum transfers
$P_tT^s_t$ in the form of ``stimulus checks'' if $z_{it}<\bar z$.

Households can save in long-term non-state-contingent nominal bonds
$\tilde b_{it}$ and capital $k_{it}$. Let $b_{it}\equiv\tilde b_{it}/P_{t-1}$ be
the real bond position at the beginning of period $t$ and $q_{t-1}$ the bond
price. Since there is no aggregate uncertainty, both assets earn the same return
along any equilibrium path except in period $0$, where the arrival of news leads
to a revaluation of asset prices and heterogeneous realized returns. Letting
$a_{it}=k_{it}+q_{t-1}b_{it}$ and $(1+r_t)$ the equalized gross return for
$t>0$:
\begin{align}
    & c_{it} + a_{it+1} \leq  w_{t} z_{it} n_{it} + (1+r_t) a_{it} - \mathcal{T}(w_{t} z_{it} n_{it}) + T_{t}^u \Ind_{z_{it}=0} + T^s_{t} \Ind_{z_{it}<\bar z}, \quad t>0 \label{budget_constraint}\\
    & c_{i0} + a_{i1} \leq  w_0 z_{i0} n_{i0} + (1+r_0^k) k_{i0} + (1+r_0^b) q_{-1} b_{i0} - \mathcal{T}(w_{0} z_{i0} n_{i0}) + T_{0}^u \Ind_{z_{i0}=0} + T^s_{0} \Ind_{z_{i0}<\bar z}
    \label{budget_constraint_t0}
\end{align}
where $w_t\equiv W_t/P_t$ is the real wage per efficiency unit and $\Ind$ is an
indicator function. I impose a standard borrowing constraint on total assets:
\begin{equation}\label{borrowing_constraint}
  a_{it+1}\geq \underline{a}.
\end{equation}

\paragraph{Lockdown constraints.} I introduce a lockdown indicator
$\text{lock}_t\in\{0,1\}$ to model the restrictions on work and consumption
observed during the pandemic. With rationing in place, the labor restriction
becomes a \emph{ceiling on the rationing index} rather than a household-level
constraint: when $\text{lock}_t=1$,
\begin{align}
& s_t \;=\; \min\!\left\{\Gamma^n,\; s^{\text{flex}}_t\right\}, \label{lockdown_n}\\
& c_{it} \;=\; \Gamma^c\, c^{ss}_i, \label{lockdown_c}
\end{align}
so hours are capped at $\Gamma^n\tilde n(z)$ and the consumption restriction
binds with equality at a common fraction $\Gamma^c$ of each household's
steady-state consumption. This setup follows \textcite{Fourakis-2024} and allows
the model to match the observed drop in output and consumption in 2020Q2.
Whether the ceiling in \eqref{lockdown_n} is the binding branch is not an
assumption but an equilibrium object, and it is verified at every solution: if
$s^{\text{flex}}_t$ fell below $\Gamma^n$ the household and government blocks
would book different hours and the solution would be invalid. Section
\ref{results} reports the realized margin for each column.

\paragraph{Recursive formulation.} Let $s\equiv\{z,a\}$ be the vector of
idiosyncratic state variables. The aggregate states are the perfect-foresight
sequences for factor prices $\{w_t\}_{t\geq0}$, $\{r_t\}_{t\geq1}$, $r^b_0$ and
$r^k_0$, labor taxes $\{\tau^n_t\}_{t\geq0}$, transfers
$\{T^s_t,T^u_t\}_{t\geq0}$, the rationing index $\{s_t\}_{t\geq0}$, lockdown
policies $\{\text{lock}_t\}_{t\geq0}$ and the transition matrices
$\{\Omega_t\}_{t\geq0}$. Note that hours no longer appear as a choice variable:
\begin{equation}
    V_{t}(z,a) = \max_{c,\,a^\prime}  \left\{(1-\beta)x^{1-1/\xi} + \beta \left[\mathbb{E}\left( V_{t+1}(z^\prime, a^\prime)^{1-\gamma}\right) \right]^{\frac{1-1/\xi}{1-\gamma}} \right\}^{\frac{1}{1-1/\xi}}
\end{equation}
subject to
\begin{align*}
        & c + a^\prime\leq  w_{t} z n + (1+r_t) a - \mathcal{T}(w_{t} z n) + T_{t}^u \Ind_{z=0} + T^s_{t}\Ind_{z<\bar z}, \quad t>0\\
    & c + a^\prime \leq  w_0 z n + (1+r_0^k) k + (1+r_0^b) q_{-1} b - \mathcal{T}(w_{0} z n) + T_{0}^u \Ind_{z=0} + T^s_{0} \Ind_{z<\bar z}\\
    & n = s_t\,\tilde n(z), \qquad a^\prime\geq \underline{a}, \qquad
      \Ind_{\text{lock}_t=1}\cdot\left(c - \Gamma^c c^{ss}\right) = 0 .
\end{align*}
Let $c^\star_{it}(z,a)$ and $a^\star_{it+1}(z,a)$ be the optimal consumption and
saving policy functions. Aggregate consumption is
$\mathcal{C}_t\equiv\int_i c^\star_{it}(z,a)\,di$, aggregate effective labor
supply is $N_t$ from \eqref{agg_labor_rationed}, and aggregate savings carried
into $t+1$ are $\mathcal{A}_{t+1}\equiv\int_i a^\star_{it+1}(z,a)\,di$.

\subsection{Foreign Households}

Foreign households demand government bonds, the only form of international
capital flows in the model. Let $B^\star_{t+1}$ denote their real position in
domestic bonds at the beginning of $t+1$, chosen in $t$:
\begin{equation} \label{foreign_demand}
B_{t+1}^\star = \Psi_t \cdot q_t^{-\chi}
\end{equation}
where $q_t$ is the bond price, $\chi>0$ governs the elasticity of foreign demand
with respect to the bond price, and $\Psi_t$ is an exogenous shifter.

\subsection{Firms}
The final good is produced by a representative firm in a perfectly competitive
market using a constant-returns-to-scale technology:
\begin{equation} \label{production_function}
    F(K_t,L_t) = Z K_t^\alpha L_t^{1-\alpha}
\end{equation}
where $K_t$ is capital, $L_t$ is effective units of labor, $\alpha$ is the
capital share and $Z$ is a productivity parameter. Firms hire labor at $w_t$,
rent capital at a pre-tax rate $\hat r^k_t$, and pay a proportional tax (or
receive a subsidy) $\tau^k_t$ on factor payments for capital. First order
conditions:
\begin{align}
    F_K(K_t,L_t) &= (1+\tau_t^k) \hat{r}_t^k \label{foc_firms1}\\
    F_L(K_t,L_t) &= w_t \label{foc_firms2}
\end{align}

% ----------------------------------------------------------------------------
\subsection{Labor Unions and Wage Setting} \label{unions}
% ----------------------------------------------------------------------------

This block is new relative to the published version of the model and is what
makes inflation endogenous.

\paragraph{Structure.} There is a continuum of labor tasks. Each task is
organized by a union that represents all households, and every household
supplies the same proportion of every task, which is what delivers the
proportional rationing in \eqref{rationing}. Union $j$ sets the nominal wage
$W_t$ for its task and faces a quadratic cost of adjusting it, as in
\textcite{Rotemberg-1982}. The structure follows
\textcite{Erceg-Henderson-Levin-2000} and its heterogeneous-agent
implementation in \textcite{Auclert-Rognlie-Straub-2024}. Unions maximize the
average utility of their members net of the adjustment cost. I express the
adjustment cost in utility units; it therefore does not appear in the goods
resource constraint \eqref{resource_constraint_open}, so measured output and
consumption are unaffected by it.

\paragraph{Wage inflation.} Define nominal wage inflation as
\begin{equation}\label{piw}
  1+\pi^w_t \;\equiv\; \frac{W_t}{W_{t-1}} \;=\; (1+\pi_t)\,\frac{w_t}{w_{t-1}},
\end{equation}
which follows from $w_t = W_t/P_t$ and $1+\pi_t = P_t/P_{t-1}$. Equation
\eqref{piw} is the bridge between the real wage the firm block prices and the
nominal rigidity the union faces: it is the only place inflation enters the real
side of the model.

\paragraph{The wage Phillips curve.} Let $\lambda_{it}$ be the marginal utility
of the GHH composite, which for the recursion \eqref{utility_function} is
\begin{equation}\label{lambda}
  \lambda_{it} \;=\; \frac{\partial V_{it}}{\partial x_{it}}
  \;=\; (1-\beta)\, V_{it}^{1/\xi}\, x_{it}^{-1/\xi},
\end{equation}
and let $\rho_{it}$ be the marginal retention rate implied by the HSV scheme
\eqref{tax_system},
\begin{equation}\label{rho}
  \rho_{it} \;=\; (1-\tau)\,\tau^n_t\,\bigl(w_t z_{it} n_{it}\bigr)^{-\tau},
\end{equation}
so that $\rho_{it} w_t z_{it}$ is the marginal after-tax real wage of household
$i$. The union's first order condition delivers the Rotemberg wage Phillips
curve
\begin{equation}\label{nkpc}
  \boxed{\;(\pi^w_t - \bar\pi)(1+\pi^w_t)
  \;=\; \frac{\kappa_w}{w_{ss}\mathcal{Y}_{ss}}\,\Xi_t
  \;+\; \beta\,(\pi^w_{t+1} - \bar\pi)(1+\pi^w_{t+1})\;}
\end{equation}
where $\bar\pi$ is the inflation target, $\kappa_w$ governs the degree of wage
stickiness, and $\Xi_t$ is the utility-weighted \emph{labor wedge}
\begin{equation}\label{wedge}
  \Xi_t \;=\; \int_i \lambda_{it}\, n_{it}
      \Bigl[\;\underbrace{\theta\, n_{it}^{1/\nu}}_{\text{marginal disutility}}
      \;-\; \underbrace{\rho_{it}\, w_t\, z_{it}}_{\text{marginal after-tax wage}}\;\Bigr] di .
\end{equation}
The wedge is the gap between the marginal rate of substitution between leisure
and consumption and the marginal after-tax real wage, weighted by each
household's marginal utility and hours. When the wedge is zero, households are
at their intratemporal optimum and wage inflation equals its target; when
households would like to work more than the union has contracted at the going
wage, the wedge is negative and wage inflation falls below target. Substituting
\eqref{rationing} into \eqref{wedge} and separating the terms that depend on the
rationing index gives the form used in computation,
\begin{equation}\label{wedge_DY}
  \Xi_t \;=\; s_t\Bigl[\theta\, s_t^{1/\nu}\,\mathcal{D}_t - w_t\,\mathcal{Y}_t\Bigr],
  \qquad
  \mathcal{D}_t \equiv \int_i \lambda_{it}\,\tilde n(z_{it})^{1+1/\nu} di,
  \qquad
  \mathcal{Y}_t \equiv \int_i \lambda_{it}\,\rho_{it}\, z_{it}\, \tilde n(z_{it})\, di,
\end{equation}
where $\mathcal{D}_t$ and $\mathcal{Y}_t$ are built from the steady-state hours
profile $\tilde n(z)$ and all of the dependence on $s_t$ is explicit. Both
integrals run over employed households only; the unemployed contribute nothing,
and $\rho_{it}$ is not defined at $z=0$. The normalization by
$w_{ss}\mathcal{Y}_{ss}$ makes $\kappa_w$ a dimensionless elasticity. In the
initial stationary equilibrium $s=1$ and $\theta=1$, so
$\Xi_{ss} = \mathcal{D}_{ss} - w_{ss}\mathcal{Y}_{ss} = 0$: the wedge vanishes
and \eqref{nkpc} reduces to $\pi^w = \bar\pi$, as it must.

% ----------------------------------------------------------------------------
\subsection{Monetary Policy} \label{monetary}
% ----------------------------------------------------------------------------

A monetary authority sets the nominal interest rate $i_t$ between $t$ and $t+1$
according to a \textcite{Taylor-1993} rule,
\begin{equation}\label{taylor}
  i_t \;=\; i^\star + \phi_\pi\,(\pi_t - \bar\pi) + \varepsilon^m_t ,
  \qquad \phi_\pi > 1,
\end{equation}
where $i^\star$ is the intercept, $\phi_\pi$ the response to inflation, and
$\varepsilon^m_t$ a monetary policy shock. The nominal rate is linked to the
real return by the Fisher equation,
\begin{equation}\label{fisher}
  1 + i_t \;=\; (1+r_{t+1})(1+\pi_{t+1}).
\end{equation}
Consistency of the rule with the terminal stationary equilibrium requires that
the intercept satisfy Fisher at that point: with $\pi=\bar\pi$,
$r=r_f$ and $\varepsilon^m\to0$,
\begin{equation}\label{istar}
  i^\star \;=\; (1+r_f)(1+\bar\pi) - 1 .
\end{equation}
This restriction matters for the counterfactuals. Any experiment that changes
the terminal real rate $r_f$ changes the right-hand side of \eqref{istar}, and
therefore changes $i^\star + \varepsilon^m_\infty$. Section
\ref{decomposition} isolates the effect of that shift.

\paragraph{Determinacy and the backward recursion.} Combining \eqref{taylor}
and \eqref{fisher} and solving for $\pi_t$ gives
\begin{equation}\label{pi_recursion}
  \pi_t \;=\; \bar\pi + \frac{(1+r_{t+1})(1+\pi_{t+1}) - 1 - i^\star - \varepsilon^m_t}{\phi_\pi}.
\end{equation}
Given the path of real returns and the terminal condition $\pi_T=\bar\pi$,
\eqref{pi_recursion} determines the entire inflation path by backward
recursion. The map contracts at rate $(1+r)/\phi_\pi \approx 0.67$ precisely
because $\phi_\pi>1$, so the Taylor principle delivers a unique bounded
solution constructively rather than merely ruling out indeterminacy. Two
features of \eqref{pi_recursion} are worth flagging because they shape the
results below. First, $i^\star$ and $\varepsilon^m_t$ enter only through their
\emph{sum}, so the intercept and the shock are not separately identified: they
are one instrument, not two. Second, the fixed point of \eqref{pi_recursion}
under a constant $C \equiv i^\star + \varepsilon^m$ is
\begin{equation}\label{fixed_point}
  \pi^\ast \;=\; \bar\pi + \frac{(1+r_f)(1+\bar\pi)-1-C}{\phi_\pi - 1 - r_f},
\end{equation}
which equals $\bar\pi$ if and only if \eqref{istar} holds. A deviation
$d$ in the rule constant therefore shifts inflation along the transition by
$-d/(\phi_\pi-1-r_f)$, an object I use directly in Section
\ref{decomposition}.

\subsection{Asset Structure}
There are two assets: a long-term nominal bond and capital. Households can
purchase a unit of the bond for a real price $q_t$. Let
$\pi_t\equiv P_t/P_{t-1}-1$, now an endogenous object. As in
\textcite{Hatchondo-Martinez-09}, a bond issued in period $t$ promises
$(\bar i+\vartheta)(1-\vartheta)^{j-1}/(1+\pi_{t+j})$ units of the final good in
period $t+j$. The constant $\bar i$ is such that the bond price equals one in
the initial pre-stimulus steady state, while $\vartheta$ controls maturity.
Bond holdings evolve as
\begin{equation}
    b_{it+1} = \ell_{it} + \left(\frac{1-\vartheta}{1+\pi_{t}}\right) b_{it},
\end{equation}
where $\ell_{it}$ is new savings. Capital depreciates at rate $\delta$.

\paragraph{Assets returns and no-arbitrage.}
\begin{enumerate}
    \item \textit{Long-term nominal bond}: a household that purchases a bond in
    $t-1$ receives a real coupon $(\bar i+\vartheta)/(1+\pi_t)$ and retains a
    fraction $(1-\vartheta)/(1+\pi_t)$ of the position, valued at $q_t$:
    \begin{equation}\label{r_b}
        1 + r_{t}^b = \frac{ (\bar{i} + \vartheta)/(1+\pi_{t}) + [(1-\vartheta)/(1+\pi_{t})] q_{t}}{q_{t-1}}
    \end{equation}
    \item \textit{Capital}: the payoff is the rental rate net of depreciation,
    \begin{equation}
        1 + r_{t}^k = 1 + \hat{r}_t^k - \delta .
    \end{equation}
\end{enumerate}
No-arbitrage equates the two for all $t>0$: $r^b_t=r^k_t=r_t$. Because the bond
is nominal and long-term, \eqref{r_b} is the channel through which endogenous
inflation redistributes: an inflation path below what was expected raises the
real value of coupon payments and transfers resources to bondholders. This
channel is absent when inflation is exogenous and fixed at $\bar\pi$, and it
turns out to account for the largest single discrepancy between the extension
and the published results (Section \ref{decomposition}).

\paragraph{Revaluation of assets.} At time $0$ realized returns are not
equalized, because households enter with positions in bonds and capital that are
revalued when news arrives. As shown in the appendix of the paper, these effects
can be represented as an additional transfer
\[
T^{\text{rev}}(q_{-1}b_{i0}, k_{i0}) \equiv (r_0^k - r_0) k_{i0} + (r_0^b - r_0) q_{-1} b_{i0},
\]
where $r_0$ is the equalized return in the initial steady state, so that
\begin{equation*}
    c_{i0} + a_{i1}\leq  w_0 z_{i0} n_{i0} + (1+r_0) a_{i0} + T^{\text{rev}}(q_{-1}b_{i0}, k_{i0})- \mathcal{T}(w_{0} z_{i0} n_{i0}) + T_{0}^u \Ind_{z=0} + T^s_{0} \Ind_{z<\bar z}.
\end{equation*}
I discipline the initial distribution of portfolios using the Survey of Consumer
Finances.

\subsection{Government}

The government raises revenue through labor and capital taxation, issues
long-term nominal bonds, and allocates resources toward interest payments,
transfers, and public consumption.

\paragraph{Taxes and transfers.}
\begin{enumerate}
    \item \textit{Labor income:} a nonlinear scheme as in
    \textcite{Heathcote-Storesletten-Violante-2017}:
    \begin{equation} \label{tax_system}
        \mathcal{T}(w_{t} z_{it} n_{it}) = w_t z_{it} n_{it} - \tau_t^n (w_t z_{it} n_{it})^{1-\tau}
    \end{equation}
    where $\tau^n_t$ controls the level of taxation and $\tau$ measures
    progressivity. Throughout, $\tau^n_t$ is a \emph{retention level}: a lower
    $\tau^n_t$ means a heavier tax burden, and $\tau^n_t>1$ corresponds to a net
    labor subsidy. This rate serves as the residual fiscal instrument that
    balances the government budget each period.
    \item \textit{Capital income:} a linear tax $\tau^k_t$ on gross capital
    returns $\hat r^k_t K_t$. In the model this becomes negative during the
    pandemic, a subsidy to firms, capturing federal support to businesses in
    2020 including tax relief and forgivable loans.
    \item \textit{Unemployment insurance:} unemployed households receive a real
    transfer $T^u_t$, independent of previous earnings.
    \item \textit{Lump-sum transfers:} real transfers $T^s_t$ conditional on
    income below an eligibility threshold, reflecting the economic impact
    payments distributed during the pandemic.
\end{enumerate}

\paragraph{Budget constraint.} With public spending $G_t$, the government
budget constraint in real terms is
\begin{equation}\label{government_budget_constraint}
    \tau_t^k \hat{r}_t^k K_t + q_t \ell_t + \int_i\mathcal{T}(w_{t} z_{it} n_{it})\, di = \left(\frac{\bar{i} + \vartheta}{1+\pi_t}\right) B_t + \int_i T_{t}^s \Ind_{z_{it}<\bar z}\, di + \int_iT_{t}^u \Ind_{z_{it}=0}\, di + G_t
\end{equation}
where $B_t$ is the stock of government bonds and
$\ell_t = B_{t+1}-[(1-\vartheta)/(1+\pi_t)]B_t$ is new issuance. Note that
inflation now enters the left-hand side of the government's problem
endogenously, through the real coupon
$(\bar i+\vartheta)/(1+\pi_t)$: the debt is nominal, so the inflation path
determines how much of the stimulus is financed by erosion of the real value of
outstanding debt rather than by taxation.

\subsection{Market Clearing}

Labor market clearing requires
\begin{equation}\label{market_clear_labor_open}
    N_t = L_t ,
\end{equation}
with $N_t$ given by \eqref{agg_labor_rationed}. Asset market clearing requires
\begin{align}
    & t > 0: \quad \mathcal{A}_{t} =  q_{t-1} \big(B_{t} - B_t^\star \big) + K_{t} \label{market_clearing_assets_open}\\
    & t = 0: \quad \mathcal{B}_0 = B_0 - B_0^\star \quad \text{and} \quad \mathcal{K}_0 = K_0 \label{market_clearing_assets_t0_open}
\end{align}
where $\mathcal{B}_0$ and $\mathcal{K}_0$ are initial aggregate holdings of
bonds and capital by domestic households and $B^\star_t$ is the foreign
position. The goods resource constraint is
\begin{equation}\label{resource_constraint_open}
    F(K_t, N_t) = \mathcal{C}_t + I_t + G_t + NX_t
\end{equation}
with $I_t\equiv K_{t+1}-(1-\delta)K_t$ and
$NX_t\equiv q_tB^\star_{t+1}-(1+r_t)q_{t-1}B^\star_t$. As noted in Section
\ref{unions}, the wage adjustment cost does not enter
\eqref{resource_constraint_open}.

% ============================================================================
\section{Equilibrium} \label{equilibrium}
% ============================================================================

\begin{shaded}
\noindent \textbf{Definition 1} (Competitive Equilibrium). \textit{
Given an initial distribution of household assets and idiosyncratic
productivities $\{b_{i0}, k_{i0}, z_{i0}\}_{i\in[0,1]}$, a sequence of
transition matrices $\{\Omega_t\}_{t\geq0}$, a government policy composed of a
fiscal policy $\{B_t,\tau^n_t,\tau^k_t,T^s_t,T^u_t,G_t\}_{t\geq0}$, a monetary
policy $\{\phi_\pi,\bar\pi,i^\star,\varepsilon^m_t\}_{t\geq0}$ and a lockdown
policy $\{\text{lock}_t\}_{t\geq0}$, an equilibrium is a sequence of aggregates
$\{\mathcal{A}_t,B^\star_t,K_t,N_t,Y_t,\mathcal{C}_t,s_t\}_{t\geq0}$ and prices
$\{q_t,r^b_t,r^k_t,w_t,\pi_t,\pi^w_t,i_t\}_{t\geq0}$ such that:
    \begin{itemize}
    \setlength\itemsep{0.4em}
        \item Given prices, transition matrices, the rationing index and
        government policy, households optimize
        \item Given prices and government policy, firms optimize
        \item \textbf{Unions optimize: the wage Phillips curve \eqref{nkpc}
        holds at every date, with the wedge \eqref{wedge} evaluated at the
        equilibrium distribution}
        \item \textbf{The monetary authority follows the Taylor rule
        \eqref{taylor}, the Fisher equation \eqref{fisher} holds, and the
        intercept satisfies \eqref{istar} at the terminal stationary
        equilibrium}
        \item \textbf{Hours are consistent with proportional rationing
        \eqref{rationing} in both the household and the government block}
        \item The sequential government budget constraint is satisfied
        \item The no-arbitrage condition is satisfied for $t>0$
        \item Markets clear
    \end{itemize}
}
\end{shaded}

\noindent A stationary equilibrium is an equilibrium in which all sequences are
constant over time. Relative to the flexible-wage version, the equilibrium
system gains one equation per date, the Phillips curve \eqref{nkpc}, and one
unknown per date, the rationing index $s_t$; inflation is then recovered from
\eqref{pi_recursion}. The bolded conditions are the new ones. The fifth
condition is not redundant: the household block and the government's tax base
must book the same hours, and it is checked at every date in the numerical
solution, since a violation would mean the labor tax revenue implied by the
household distribution differs from what the budget constraint records.

% ============================================================================
\section{Calibration} \label{calibration}
% ============================================================================

The model is calibrated to the U.S. economy using quarterly data spanning 2019
to 2024. The initial stationary equilibrium is calibrated using data from 2019Q1
to 2019Q4. The transition path begins in 2020Q2, coinciding with the onset of
the pandemic and the expansion of federal government debt, and I use observed
time series through 2024Q4 to discipline the evolution of key aggregates.
Table \ref{tab:calib} lists the parameters, separating those inherited from the
published calibration from the three that the sticky-wage block adds.

\begin{table}[t]\centering
\caption{Calibration}\label{tab:calib}
\small
\begin{tabular}{llr}
\toprule
Parameter & Description & Value \\
\midrule
\multicolumn{3}{l}{\textit{Preferences and technology (unchanged)}}\\
$\beta$ & Discount factor & $0.988043$ \\
$\xi$ & Elasticity of intertemporal substitution & $1.00$ \\
$\gamma$ & Risk aversion & $4.00$ \\
$\nu$ & Frisch elasticity of labor supply & $1.00$ \\
$\theta$ & Labor disutility & $1.00$ \\
$\alpha$ & Capital share & $0.33$ \\
$\delta$ & Depreciation (quarterly) & $0.014292$ \\
$Z$ & Productivity & $0.105810$ \\
$\underline{a}$ & Borrowing limit & $-0.02$ \\
$\chi$ & Elasticity of foreign bond demand & $0.44$ \\
$\tau$ & HSV progressivity & $0.181$ \\
$P^{ee}$, $P^{uu}$ & Steady-state labor market transitions & $0.975$, $0.400$ \\
\midrule
\multicolumn{3}{l}{\textit{Sticky-wage and monetary block (new)}}\\
$\kappa_w$ & Slope of the wage Phillips curve & $0.05$ \\
$\phi_\pi$ & Taylor rule response to inflation & $1.50$ \\
$\bar\pi$ & Inflation target (annualized) & $2.00\%$ \\
$i^\star$ & Rule intercept, from \eqref{istar} (annualized) & $3.0994\%$ \\
\midrule
\multicolumn{3}{l}{\textit{Lockdown, 2020Q2 (unchanged targets)}}\\
$\Gamma^n$ & Hours ceiling, fraction of steady-state hours & $0.902301$ \\
$\Gamma^c$ & Consumption restriction, fraction of steady state & $0.922960$ \\
\midrule
\multicolumn{3}{l}{\textit{Initial stationary equilibrium (implied)}}\\
$r_{ss}$ & Real rate (annualized) & $0.9842\%$ \\
$w_{ss}$ & Real wage & $0.095539$ \\
$\tau^n_{ss}$ & HSV retention level & $0.886550$ \\
$\tau^k_{ss}$ & Capital income tax & $0.137931$ \\
$r_f$ & Terminal real rate, baseline (annualized) & $1.0778\%$ \\
\bottomrule
\end{tabular}
\end{table}

The monetary shock path $\{\varepsilon^m_t\}$ is not a free parameter. It is
recovered as the residual of the Taylor rule \eqref{taylor} evaluated at the
solved real rate path and the \emph{observed} inflation series, so that the
baseline transition reproduces measured inflation by construction:
\begin{equation}\label{epsm}
  \varepsilon^m_t \;=\; \bigl[(1+r_{t+1})(1+\pi^{\text{data}}_{t+1}) - 1\bigr]
  \;-\; i^\star \;-\; \phi_\pi\bigl(\pi^{\text{data}}_t - \bar\pi\bigr).
\end{equation}
This is the natural analogue of the exogenous inflation assumption in the
published model: rather than imposing the inflation path directly, I impose the
monetary shocks that generate it, which lets the model speak to how much of
observed inflation the fiscal and foreign blocks can account for on their own.
The answer, reported in Section \ref{results}, is very little.

% ============================================================================
\section{Counterfactual Results} \label{results}
% ============================================================================

I consider four counterfactuals, each holding the monetary rule and the monetary
shocks fixed and changing one element of the fiscal or external environment.
Welfare is measured as a consumption-equivalent (CE) gain at the announcement
date relative to the pre-stimulus stationary equilibrium.

\paragraph{The monetary rule is held fixed across columns.} Because $i^\star$
and $\varepsilon^m_t$ enter \eqref{pi_recursion} only as their sum, the object
that must be held fixed is the rule constant
$C_t \equiv i^\star + \varepsilon^m_t$. Two of the four counterfactuals move the
terminal real rate $r_f$ and therefore, by \eqref{istar}, the value the constant
must take in the far tail. I hold $C_t$ at its baseline value for
$t \leq 600$ quarters and let it revert to the counterfactual's own
Fisher-consistent value thereafter. Since $\beta^{600}=7.3\times10^{-4}$, this
has no effect on announcement-date welfare, while keeping the terminal condition
$\pi_T=\bar\pi$ exactly consistent. Under this specification the rule constant
is numerically identical across all four columns at every date
$t=2,\dots,600$. Table \ref{tab:rule} records the rule for each column.

\begin{table}[t]\centering
\caption{The monetary rule by column}\label{tab:rule}
\small
\begin{tabular}{lrrrr}
\toprule
Column & $i^\star$ (ann.) & $\phi_\pi$ & $r_f$ (ann.) & Constant adjustment \\
\midrule
Baseline                 & $3.0994\%$ & $1.50$ & $1.0778\%$ & none \\
No capital subsidies     & $3.0994\%$ & $1.50$ & $1.0778\%$ & none \\
Lower foreign demand     & $3.1531\%$ & $1.50$ & $1.1305\%$ & $t>600$ only \\
Only stimulus checks     & $3.0994\%$ & $1.50$ & $1.0778\%$ & none \\
CBO debt projections     & $3.4495\%$ & $1.50$ & $1.4211\%$ & $t>600$ only \\
\bottomrule
\end{tabular}
\end{table}

\paragraph{Welfare groups.} Groups are defined by quantiles of the asset grid
weighted by the invariant wealth distribution of the initial stationary
equilibrium. Because a mass of households sits exactly at the borrowing
constraint, the realized masses are not exact deciles: the bottom bin holds
$15.3\%$ of households rather than $10\%$, and the bins above it are each
shifted up by roughly five percentage points. ``Poor'' is therefore the bottom
$15.3\%$, ``Middle'' spans about the $44$th to $63$rd percentile, and ``Rich''
is the top $9.3\%$. I report the labels of the published table for
comparability, with this correction noted.

\begin{table}[t]\centering
\caption{Consumption-equivalent welfare gains (\%)}\label{tab:table6}
\small
\begin{tabular}{lrrrrr}
\toprule
 & Baseline & No capital & Lower foreign & Only stimulus & CBO debt \\
 &          & subsidies  & demand        & checks        & projections \\
\midrule
\multicolumn{6}{l}{\textit{(a) Levels, against the pre-stimulus steady state}}\\
Overall mean       & $1.910$ & $2.106$ & $1.236$  & $4.449$ & $2.245$ \\
Overall minimum    & $-5.383$& $-6.096$& $-10.063$& $-7.251$& $-3.092$ \\
Poor (bottom 10\%) & $3.350$ & $4.085$ & $2.229$  & $8.929$ & $3.844$ \\
Middle (40--60\%)  & $2.043$ & $2.314$ & $1.350$  & $4.422$ & $2.294$ \\
Rich (top 10\%)    & $-0.325$& $-1.848$& $-0.215$ & $-2.037$& $0.369$ \\
\midrule
\multicolumn{6}{l}{\textit{(b) Difference from the baseline (pp)}}\\
Overall mean       & --- & $+0.196$ & $-0.674$ & $+2.539$ & $+0.334$ \\
Overall minimum    & --- & $-0.713$ & $-4.680$ & $-1.868$ & $+2.291$ \\
Poor (bottom 10\%) & --- & $+0.735$ & $-1.121$ & $+5.579$ & $+0.494$ \\
Middle (40--60\%)  & --- & $+0.271$ & $-0.693$ & $+2.378$ & $+0.251$ \\
Rich (top 10\%)    & --- & $-1.523$ & $+0.110$ & $-1.712$ & $+0.694$ \\
\midrule
\multicolumn{6}{l}{\textit{(c) Published values, flexible-wage model}}\\
Overall mean       & $1.82$ & $2.02$ & $1.22$ & $4.29$ & $2.51$ \\
Overall minimum    & $-1.77$& $-6.18$& $-3.72$& $-5.80$& $-1.27$ \\
Poor (bottom 10\%) & $3.03$ & $3.77$ & $2.03$ & $8.60$ & $4.00$ \\
Middle (40--60\%)  & $2.00$ & $2.27$ & $1.34$ & $4.28$ & $2.62$ \\
Rich (top 10\%)    & $-0.20$& $-1.69$& $0.06$ & $-1.90$& $0.50$ \\
\bottomrule
\end{tabular}
\end{table}

\paragraph{Reading Table \ref{tab:table6}.} On every average, the extension
reproduces the published levels closely: the overall mean is within $0.02$ to
$0.27$ percentage points across all five columns, Middle within $0.33$, Rich
within $0.28$, and Poor within $0.33$. The ranking and sign pattern of the
published table survive: removing the capital subsidies raises average welfare
while making the rich substantially worse off; a weaker foreign demand for
government debt lowers welfare throughout; a stimulus package consisting of
checks alone is far more progressive; and matching the CBO debt path raises
welfare relative to the realized path.

The exception is the \textit{Overall minimum} row, which I do not regard as
reconciled. That statistic is the CE gain of a single $(a,z)$ cell with positive
mass, so it is far more sensitive to the specification than any average. The
extension puts it at $-5.383$ in the baseline against $-1.77$ in the published
model, and at $-10.063$ against $-3.72$ under lower foreign demand, while
matching to $0.08$ under no capital subsidies. The direction is intuitive
--- with sticky wages, the households worst hit by the rationing of hours have no
margin on which to adjust --- but I have not isolated the mechanism, and the row
should be read as an artifact of the extension rather than a result.

\paragraph{What the extension adds.} Table \ref{tab:extension} reports the
objects the flexible-wage model cannot produce. The interpretation of the last
column requires care. Because $\varepsilon^m_t$ is recovered from \eqref{epsm}
so that the baseline matches measured inflation, the peak deviation of about
$6.8$ percentage points is the deviation \emph{in the data}, not an endogenous
model response. What the extension generates on its own is the difference across
columns, and that is small: at most $+0.15$ percentage points at 2022Q2, under
lower foreign demand. Absent monetary shocks, the fiscal and external blocks of
this model move inflation very little --- which is the central quantitative
finding of the extension.

\begin{table}[t]\centering
\caption{Inflation and real rates under the extension}\label{tab:extension}
\small
\begin{tabular}{lrrrrr}
\toprule
Column & $\pi$ 2022Q2 & vs base & $r$ at $t=20$ & vs base & Peak $\pi$ dev. \\
       & (ann.\ \%)   & (pp)    & (ann.\ \%)    & (pp)    & from target (pp) \\
\midrule
Baseline              & $9.0600$ & ---       & $0.9862$ & ---       & $+6.7823$ \\
No capital subsidies  & $9.0624$ & $+0.0022$ & $0.9878$ & $+0.0016$ & $+6.7845$ \\
Lower foreign demand  & $9.2212$ & $+0.1510$ & $1.0530$ & $+0.0663$ & $+6.9333$ \\
Only stimulus checks  & $9.0679$ & $+0.0074$ & $0.9928$ & $+0.0065$ & $+6.7898$ \\
CBO debt projections  & $9.1248$ & $+0.0607$ & $1.1086$ & $+0.1215$ & $+6.8430$ \\
\bottomrule
\end{tabular}
\\[0.5em]
\parbox{0.93\textwidth}{\footnotesize The peak deviation occurs at 2022Q2 in
every column. Terminal inflation returns to the $2\%$ target in all five runs.
The labor tax retention level $\tau^n$ at the announcement date is $0.8919$ in
the baseline, $1.3583$ under no capital subsidies (a net labor subsidy),
$0.8283$ under lower foreign demand, $0.8950$ under only stimulus checks and
$0.9456$ under the CBO debt path.}
\end{table}

\paragraph{Solution diagnostics.} Table \ref{tab:diag} records the numerical
status of each column. The item labelled ``hours gate'' is the consistency check
required by the fifth condition of Definition 1: the labor tax revenue implied
by integrating the household distribution against the HSV schedule, compared
with what the government budget constraint records. It holds to a relative
$1.6\times10^{-7}$ in every column, which is the mass normalization error of the
distribution rather than an error in the block. The ``lockdown margin'' column
reports $s^{\text{flex}}_0/\Gamma^n - 1$: it is positive everywhere, so the
ceiling in \eqref{lockdown_n} is the binding branch at the announcement date in
every column, as required. All five columns deliver the same $8.4002\%$ fall in
2020Q2 output, since $s_0=\Gamma^n$ by construction.

\begin{table}[t]\centering
\caption{Solution diagnostics}\label{tab:diag}
\small
\begin{tabular}{lrrrrr}
\toprule
Column & Iterations & Wall (min) & Residual & Hours gate & Lockdown margin \\
\midrule
Baseline              & $1$  & $0.19$  & $7.9\times10^{-6}$ & $1.6\times10^{-7}$ & $+14.78\%$ \\
No capital subsidies  & $37$ & $6.80$  & $4.7\times10^{-6}$ & $1.6\times10^{-7}$ & $+63.90\%$ \\
Lower foreign demand  & $44$ & $8.22$  & $8.8\times10^{-6}$ & $1.6\times10^{-7}$ & $+7.82\%$ \\
Only stimulus checks  & $40$ & $6.62$  & $9.5\times10^{-6}$ & $1.6\times10^{-7}$ & $+15.12\%$ \\
CBO debt projections  & $86$ & $16.11$ & $9.5\times10^{-6}$ & $1.6\times10^{-7}$ & $+20.61\%$ \\
\bottomrule
\end{tabular}
\end{table}

\paragraph{Distributional detail.} Table \ref{tab:deciles} gives the full
profile. The pattern that distinguishes the columns is the behavior of the top
bin. Removing the capital subsidies and moving to checks alone both impose large
losses on the top decile ($-1.52$ and $-1.71$ percentage points relative to the
baseline) while benefiting every other bin. The CBO debt path is the only
counterfactual that raises welfare at the top, and it does so by more than it
raises welfare in the middle --- a pattern that the next section attributes to a
specific channel.

\begin{table}[t]\centering
\caption{Difference from baseline by wealth bin (pp)}\label{tab:deciles}
\small
\begin{tabular}{lrrrrr}
\toprule
Bin & Realized mass & No capital & Lower foreign & Only stimulus & CBO debt \\
\midrule
1 (poorest) & $15.26\%$ & $+0.7355$ & $-1.1208$ & $+5.5787$ & $+0.4940$ \\
2           & $10.16\%$ & $+0.7042$ & $-1.0562$ & $+4.6131$ & $+0.4555$ \\
3           & $9.53\%$  & $+0.5124$ & $-0.9303$ & $+3.9600$ & $+0.4334$ \\
4           & $9.27\%$  & $+0.4122$ & $-0.8219$ & $+3.2314$ & $+0.3855$ \\
5           & $9.08\%$  & $+0.3064$ & $-0.7371$ & $+2.6653$ & $+0.2965$ \\
6           & $9.25\%$  & $+0.2350$ & $-0.6484$ & $+2.0915$ & $+0.2046$ \\
7           & $9.40\%$  & $+0.1353$ & $-0.5404$ & $+1.5240$ & $+0.1186$ \\
8           & $9.31\%$  & $+0.0674$ & $-0.4235$ & $+0.9656$ & $+0.0571$ \\
9           & $9.39\%$  & $-0.0184$ & $-0.2568$ & $+0.3538$ & $+0.0920$ \\
10 (richest)& $9.33\%$  & $-1.5233$ & $+0.1101$ & $-1.7123$ & $+0.6940$ \\
\bottomrule
\end{tabular}
\end{table}

% ============================================================================
\section{The Neutral-Rate Channel} \label{decomposition}
% ============================================================================

The two counterfactuals that change the terminal real rate --- lower foreign
demand and the CBO debt path --- activate a channel that does not exist in the
flexible-wage model, and it is worth isolating because it accounts for the
largest discrepancy between the extension and the published results.

\paragraph{The channel.} By \eqref{istar}, a counterfactual that raises the
terminal real rate $r_f$ raises the Fisher-consistent rule constant. By
\eqref{fixed_point}, holding the monetary shocks fixed, a rise $d$ in that
constant lowers inflation along the transition by
\begin{equation}\label{intercept_channel}
  d\pi \;=\; -\,\frac{d}{\phi_\pi - 1 - r_f}.
\end{equation}
Since government debt is nominal and long-term, a permanently lower inflation
path raises the real value of coupon payments, transferring resources to
bondholders, while the heavier labor taxation needed to service the debt falls
on households who hold little of it. The channel is thus redistributive by
construction, and it operates in the direction of the wealthy.

\paragraph{Measurement.} I switch the channel off by subtracting $d$ from the
monetary shock over the transition, which restores the rule constant to its
baseline value while leaving the terminal condition intact --- the post-600
specification of Section \ref{results}. Comparing that solution with one in
which the constant is allowed to rise isolates the channel. The measured shift
in inflation matches \eqref{intercept_channel} to $0.0008$ percentage points for
the CBO column and $0.0001$ for lower foreign demand, flat across four hundred
quarters, while real rates move by at most $0.006$ percentage points. The
channel is, to a very good approximation, purely nominal.

\begin{table}[t]\centering
\caption{Contribution of the neutral-rate channel to welfare (pp)}\label{tab:channel}
\small
\begin{tabular}{lrr}
\toprule
Bin & CBO debt projections & Lower foreign demand \\
\midrule
1 (poorest) & $+0.449$ & $+0.069$ \\
2           & $+0.512$ & $+0.078$ \\
3           & $+0.709$ & $+0.108$ \\
4           & $+0.842$ & $+0.128$ \\
5           & $+0.801$ & $+0.122$ \\
6           & $+0.731$ & $+0.111$ \\
7           & $+0.626$ & $+0.095$ \\
8           & $+0.451$ & $+0.068$ \\
9           & $+0.227$ & $+0.034$ \\
10 (richest)& $-0.914$ & $-0.137$ \\
\midrule
Overall mean & $+0.443$ & $+0.068$ \\
Poor         & $+0.449$ & $+0.069$ \\
Middle       & $+0.766$ & $+0.116$ \\
Rich         & $-0.914$ & $-0.137$ \\
\bottomrule
\end{tabular}
\\[0.5em]
\parbox{0.83\textwidth}{\footnotesize Entries are the change in the CE gain from
switching the channel off. A positive entry means the channel was reducing that
bin's welfare.}
\end{table}

\paragraph{Result.} For the CBO column the channel is large and sharply
redistributive. With it operating, the top bin gains $+1.609$ percentage points
relative to the baseline; switching it off removes $0.914$ of that, so
\textbf{the channel accounts for $57\%$ of the entire gain accruing to the top
decile}. Over bins three through eight it works in the opposite direction, and
removing it converts losses of $0.28$--$0.53$ percentage points into gains of
$0.06$--$0.43$. Without the channel, all four welfare groups gain in the CBO
column, as they do in the published table, and the top group matches the
published value to $0.006$ percentage points.

For lower foreign demand the channel is proportionally similar but small in
absolute terms, at most $0.14$ percentage points against a counterfactual effect
of $-1.19$ to $+0.25$: the terminal rate moves by only $5$ basis points there,
against $34$ for the CBO path. Indeed the channel's welfare contribution scales
almost exactly linearly in $d$: dividing the CBO contribution by the ratio of
the two shifts, $6.51$, predicts the lower-foreign contribution to within one to
three percent in every group. That linearity is useful, since it means the
channel can be priced for any further experiment that moves the terminal steady
state from the size of the shift alone.

\paragraph{Interpretation.} The published, flexible-wage model has no such
channel: with inflation exogenous and constant, a change in the terminal real
rate cannot redistribute through the nominal value of government debt. The
extension therefore does not overturn the paper's conclusions, but it does
identify an additional force that operates whenever a fiscal experiment moves
the neutral rate, and it shows that force to be quantitatively first-order for
the debt-path counterfactual specifically. A reader of the published table
should treat the CBO column's distributional pattern as the one most sensitive
to the assumption that inflation is exogenous.

\end{document}
